\clearpage
\thispagestyle{empty}
\vspace*{\fill}

\begin{minipage}{0.88\textwidth}
\raggedright
{\zihao{3}\bfseries 版权说明\par}
\vspace{0.8em}

\renewcommand{\arraystretch}{1.25}
\begin{tabular}{@{}p{0.18\textwidth}p{0.64\textwidth}@{}}
作品名称 & 《一个月打通高中物理》\\
副标题 & 写给初三升高一学生和竞赛入门者的物理预习讲义\\
著者 & 董兴习、许琴\\
版本 & 电子讲义修订版\\
版本日期 & 2026 年 7 月\\
\end{tabular}

\vspace{1.2em}

\textcopyright\ 2026 董兴习、许琴。保留所有依法享有的权利。

本书中具有独创性的文字表达、内容编排、例题设计和自绘图示，其相关权利依法归著作权人所有。物理学中的通用概念、定律、公式、符号和公共知识不因收录于本书而成为排他性内容。

未经著作权人书面许可，不得以复制、发行、信息网络传播、改编、翻译、汇编等方式使用本书全部或实质性部分，但法律法规另有规定的除外。个人学习者可以打印一份供本人非商业学习使用；课堂教学中依法进行少量引用或复制时，应注明作品名称和著者，不得用于出版发行或商业销售。

\textbf{内容说明：} 本书是面向初三升高一学生和高中物理竞赛入门者的辅助学习材料，不是教育行政部门审定教材，也不能替代课程标准、学校教学安排和教师的专业指导。不同地区的课程进度和考试要求可能不同，使用时应以当地现行要求为准。

\textbf{版本说明：} 本文件为持续校订的电子作品版本。若以后由出版机构正式出版，书名、版次、出版信息和版权页内容以正式出版版本为准。

\textbf{读者反馈：} 如发现概念、公式、例题解答、图示或排版问题，请通过“怎样使用这本讲义”末尾的联系方式向作者反馈。
\end{minipage}

\vspace*{0.08\textheight}
\clearpage
